% Laget automatisk av elevark/lag_elevark.py fra opplegg/ballistisk-pendel.html. Endringer her blir overskrevet.
\documentclass[11pt]{article}
\usepackage{elevark}
\begin{document}
\elevarkhode{Fysikk 1 · Bevegelsesmengde og støt}{Ballistisk pendel}{https://fysikklab.no/opplegg/ballistisk-pendel.html}{Opplegget på nett}

En kule skytes inn i en pendel og blir sittende fast. Hvor fort gikk kula? Dere skal finne farten på to uavhengige måter: fra hvor langt kula flyr når den skytes horisontalt ut over gulvet, og fra hvor høyt pendelen svinger. Gir metodene samme svar?

\section*{Hensikt}

\begin{itemize}
\item Bestemme farten til en kule med horisontalt kast og med bevaring av bevegelsesmengde og energi.
\item Vurdere om metodene stemmer innenfor usikkerheten.
\item Finne ut hvor stor del av den kinetiske energien som går tapt i et fullstendig uelastisk støt.
\end{itemize}

\section*{Utstyr}

\begin{itemize}
\item Ballistisk pendel med fjærkanon og kule, og en vinkelviser som blir stående i det største utslaget.
\item Målebånd, karbonpapir (eller et ark med kopipapir over) på gulvet, og en vekt.
\item Vernebriller. Ingen står foran kanonen når den er spent.
\end{itemize}

\section*{Teorien}

\textbf{Kastet.} Kula forlater kanonen horisontalt i høyden $h$ over gulvet. Falltiden er $t = \sqrt{2h/g}$, og farten er $v = x/t = x \cdot \sqrt{g/(2h)}$, der $x$ er kastlengden.

\textbf{Pendelen.} I støtet er bevegelsesmengden bevart: $mv = (m + M)V$, der $m$ er massen til kula, $M$ massen til pendelen og $V$ farten rett etter støtet. Etterpå svinger pendelen opp, og mekanisk energi er bevart: $\tfrac{1}{2}(m + M)V^{2} = (m + M)g \cdot L(1 - \cos θ)$. Her er $L$ avstanden fra opphenget til tyngdepunktet til pendelen med kula i, og $θ$ er det største utslaget. Det gir $V = \sqrt{2gL(1 - \cos  θ)}$ og $v = (m + M)/m \cdot V$.

\section*{Framgangsmåte}

\begin{enumerate}
\item Vei kula og pendelen (uten kula). Sett kula inn i pendelen og finn tyngdepunktet ved å balansere pendelen på en kant. Mål $L$ fra aksen til tyngdepunktet.
\item \textbf{Kastet:} Vipp pendelen til side. Mål høyden $h$ fra underkanten av kula i kanonen til gulvet. Skyt én gang for å se hvor kula lander, og legg karbonpapiret der. Skyt fem ganger på samme spenning, og mål kastlengden $x$ fra rett under kanonmunningen til hvert merke.
\item \textbf{Pendelen:} Heng pendelen ned igjen, sett vinkelviseren på null og skyt fem ganger på samme spenning. Les av det største utslaget $θ$ hver gang.
\item Skriv målingene inn i tabellen. Hver rad er ett skudd med hver metode.
\end{enumerate}

\section*{Oppgaver}

\begin{oppgaveliste}
\oppgave \fase{gjett}{Gjett}Hvor stor del av den kinetiske energien til kula tror dere er igjen rett etter at den har satt seg fast i pendelen?
\svar
\oppgave \fase{test}{Mål}Gjør målingene, eller trykk på «Fyll inn eksempel». Hvilken fart får dere med hver metode? Stemmer metodene innenfor usikkerheten?
\svar
\oppgave \fase{test}{Regn}Vis at den kinetiske energien rett etter støtet er $m/(m + M)$ av energien før støtet. Hvor mye energi går tapt i eksemplet?
\svar
\oppgave \fase{forklar}{Forklar}Hvorfor kan vi bruke bevaring av bevegelsesmengde i selve støtet, men ikke bevaring av mekanisk energi? Hvorfor kan vi bruke energibevaring når pendelen svinger opp?
\svar
\oppgave \fase{forklar}{Vurder}Hvor mye endres farten fra pendelen hvis dere leser av $θ$ én grad for stort? Hvor mye endres farten fra kastet hvis $x$ er 1~cm for langt? Hvilken måling er viktigst å gjøre nøyaktig?
\svar
\oppgave \fase{utfordring}{Utfordring}Tenk dere at kula i stedet sprettet rett tilbake fra pendelen. Ville pendelen svingt høyere eller lavere enn når kula setter seg fast? Forklar med bevegelsesmengde.
\svar
\end{oppgaveliste}

\end{document}
